\documentclass[11pt,a4paper]{article}

\usepackage[utf8]{inputenc}
\usepackage[T1]{fontenc}
\usepackage{amsmath,amssymb,mathtools}
\usepackage{xcolor}
\usepackage{mdframed}
\usepackage{geometry}
\usepackage{microtype}

\geometry{
  top=2.3cm,
  bottom=2.3cm,
  left=2.4cm,
  right=2.4cm
}

\setlength{\parindent}{0pt}
\setlength{\parskip}{0.7em}

%------------------------------------------------
% Couleurs
%------------------------------------------------

\definecolor{bleutheoreme}{RGB}{31,78,121}
\definecolor{fondbleu}{RGB}{235,243,250}

\definecolor{orangelemme}{RGB}{210,105,30}
\definecolor{fondorange}{RGB}{255,244,230}

\definecolor{vertpreuve}{RGB}{35,130,75}
\definecolor{fondvert}{RGB}{235,248,239}

\definecolor{violetremarque}{RGB}{112,48,160}
\definecolor{fondviolet}{RGB}{246,238,252}

\definecolor{rougeresultat}{RGB}{180,30,45}
\definecolor{fondrouge}{RGB}{255,235,238}

%------------------------------------------------
% Environnements encadr\'es
%------------------------------------------------

\newmdenv[
  linecolor=bleutheoreme,
  backgroundcolor=fondbleu,
  linewidth=1.2pt,
  roundcorner=4pt,
  skipabove=10pt,
  skipbelow=10pt,
  innerleftmargin=10pt,
  innerrightmargin=10pt
]{theoremebox}

\newmdenv[
  linecolor=orangelemme,
  backgroundcolor=fondorange,
  linewidth=1.2pt,
  roundcorner=4pt,
  skipabove=10pt,
  skipbelow=10pt,
  innerleftmargin=10pt,
  innerrightmargin=10pt
]{lemmebox}

\newmdenv[
  linecolor=vertpreuve,
  backgroundcolor=fondvert,
  linewidth=1.2pt,
  roundcorner=4pt,
  skipabove=10pt,
  skipbelow=10pt,
  innerleftmargin=10pt,
  innerrightmargin=10pt
]{preuvebox}

\newmdenv[
  linecolor=violetremarque,
  backgroundcolor=fondviolet,
  linewidth=1.2pt,
  roundcorner=4pt,
  skipabove=10pt,
  skipbelow=10pt,
  innerleftmargin=10pt,
  innerrightmargin=10pt
]{remarquebox}

\newmdenv[
  linecolor=rougeresultat,
  backgroundcolor=fondrouge,
  linewidth=1.5pt,
  roundcorner=4pt,
  skipabove=12pt,
  skipbelow=12pt,
  innerleftmargin=10pt,
  innerrightmargin=10pt
]{resultatbox}

%------------------------------------------------
% Commandes math\'ematiques
%------------------------------------------------

\newcommand{\N}{\mathbb{N}}
\newcommand{\R}{\mathbb{R}}
\newcommand{\1}{\mathbf{1}}

%------------------------------------------------
% Document
%------------------------------------------------

\title{
  \textbf{Somme de puissances \`a exposant croissant}\\
  \large D\'emonstration par convergence domin\'ee
}
\author{R\'edaction rigoureuse pour Mostafa}
\date{12 juillet 2026}

\begin{document}

\maketitle

\begin{theoremebox}
\textbf{\'Enonc\'e.}

Soit \(\alpha>0\). D\'emontrer que, lorsque
\(n\longrightarrow+\infty\),
\[
1^{\alpha n}+2^{\alpha n}+\cdots+n^{\alpha n}
\sim
\frac{n^{\alpha n}}{1-e^{-\alpha}}.
\]
\end{theoremebox}

\begin{lemmebox}
\textbf{Lemme de domination.}

Pour tous entiers \(n\geq 1\) et \(k\in\N\), posons
\[
f_n(k)
=
\1_{\{0,\ldots,n-1\}}(k)
\left(1-\frac{k}{n}\right)^{\alpha n}.
\]
Alors
\[
0\leq f_n(k)\leq e^{-\alpha k}.
\]
\end{lemmebox}

\begin{preuvebox}
\textbf{Preuve du lemme.}

Soient \(n\geq1\) et \(k\in\N\).

Si \(k\geq n\), alors
\[
f_n(k)=0,
\]
et l'in\'egalit\'e
\[
0\leq f_n(k)\leq e^{-\alpha k}
\]
est satisfaite.

Supposons maintenant que \(0\leq k\leq n-1\). La fonction
\[
h:[0,1)\longrightarrow\R,
\qquad
h(x)=\log(1-x)+x,
\]
est d\'erivable et v\'erifie
\[
h'(x)
=
-\frac{1}{1-x}+1
=
-\frac{x}{1-x}
\leq 0.
\]
Elle est donc d\'ecroissante sur \([0,1)\). Comme \(h(0)=0\),
on obtient, pour tout \(x\in[0,1)\),
\[
\log(1-x)\leq -x.
\]

En appliquant cette in\'egalit\'e \`a \(x=k/n\), il vient
\[
\log\left(1-\frac{k}{n}\right)
\leq
-\frac{k}{n}.
\]
Puisque \(\alpha n>0\),
\[
\alpha n
\log\left(1-\frac{k}{n}\right)
\leq
-\alpha k.
\]
La fonction exponentielle \'etant croissante sur \(\R\), on en
d\'eduit
\[
\left(1-\frac{k}{n}\right)^{\alpha n}
\leq e^{-\alpha k}.
\]
Dans le cas \(0\leq k\leq n-1\), le membre de gauche est
positif. Par cons\'equent,
\[
0\leq f_n(k)\leq e^{-\alpha k}.
\]
\end{preuvebox}

\begin{preuvebox}
\textbf{Preuve principale par convergence domin\'ee.}

Posons
\[
S_n=\sum_{j=1}^{n}j^{\alpha n}.
\]
Apr\`es division par \(n^{\alpha n}\), nous obtenons
\[
\frac{S_n}{n^{\alpha n}}
=
\sum_{j=1}^{n}
\left(\frac{j}{n}\right)^{\alpha n}.
\]

Effectuons le changement d'indice
\[
k=n-j.
\]
Lorsque \(j\) parcourt \(\{1,\ldots,n\}\), l'entier \(k\)
parcourt \(\{0,\ldots,n-1\}\). Il vient donc
\[
\frac{S_n}{n^{\alpha n}}
=
\sum_{k=0}^{n-1}
\left(1-\frac{k}{n}\right)^{\alpha n}.
\]

Consid\'erons l'espace mesur\'e
\[
\bigl(\N,\mathcal{P}(\N),\mu\bigr),
\]
o\`u \(\mu\) d\'esigne la mesure de comptage. Pour toute fonction
positive \(u:\N\to[0,+\infty]\), on a
\[
\int_{\N}u\,d\mu
=
\sum_{k=0}^{\infty}u(k).
\]

Pour \(n\geq1\), d\'efinissons
\[
f_n:\N\longrightarrow[0,+\infty)
\]
par
\[
f_n(k)
=
\begin{cases}
\displaystyle
\left(1-\frac{k}{n}\right)^{\alpha n},
& 0\leq k\leq n-1,\\[1.2ex]
0,
& k\geq n.
\end{cases}
\]
On a alors
\[
\frac{S_n}{n^{\alpha n}}
=
\sum_{k=0}^{\infty}f_n(k)
=
\int_{\N}f_n\,d\mu.
\]

Nous v\'erifions successivement la convergence simple et
l'existence d'une fonction dominante int\'egrable.

\medskip

\textbf{1. Convergence simple.}

Fixons \(k\in\N\).

Si \(k=0\), alors, pour tout \(n\geq1\),
\[
f_n(0)=1.
\]
Par cons\'equent,
\[
\lim_{n\to\infty}f_n(0)=1=e^{-\alpha\cdot0}.
\]

Supposons \(k\geq1\). Pour tout \(n>k\), on a
\[
f_n(k)
=
\left(1-\frac{k}{n}\right)^{\alpha n}.
\]
En passant au logarithme,
\[
\log f_n(k)
=
\alpha n
\log\left(1-\frac{k}{n}\right).
\]
Nous pouvons \'ecrire
\[
n\log\left(1-\frac{k}{n}\right)
=
-k\,
\frac{
  \log\left(1-\frac{k}{n}\right)
}{
  -\frac{k}{n}
}.
\]
Or
\[
\frac{k}{n}\longrightarrow0
\]
et
\[
\lim_{x\to0}
\frac{\log(1-x)}{-x}
=
1.
\]
Ainsi,
\[
n\log\left(1-\frac{k}{n}\right)
\longrightarrow -k.
\]
Apr\`es multiplication par \(\alpha\),
\[
\log f_n(k)\longrightarrow-\alpha k.
\]
Par continuit\'e de la fonction exponentielle,
\[
f_n(k)\longrightarrow e^{-\alpha k}.
\]

Nous avons donc, pour tout \(k\in\N\),
\[
\lim_{n\to\infty}f_n(k)=e^{-\alpha k}.
\]

\medskip

\textbf{2. Domination par une fonction int\'egrable.}

D'apr\`es le lemme, pour tous \(n\geq1\) et \(k\in\N\),
\[
0\leq f_n(k)\leq e^{-\alpha k}.
\]
Posons
\[
g(k)=e^{-\alpha k}.
\]
La fonction \(g\) est int\'egrable pour la mesure de comptage si
et seulement si la s\'erie
\[
\sum_{k=0}^{\infty}e^{-\alpha k}
\]
converge.

Comme \(\alpha>0\), on a
\[
0<e^{-\alpha}<1.
\]
La s\'erie pr\'ec\'edente est donc une s\'erie g\'eom\'etrique
convergente, de somme
\[
\sum_{k=0}^{\infty}e^{-\alpha k}
=
\frac{1}{1-e^{-\alpha}}.
\]
En particulier,
\[
\int_{\N}g\,d\mu
=
\frac{1}{1-e^{-\alpha}}
<+\infty.
\]

\medskip

\textbf{3. Application du th\'eor\`eme de convergence domin\'ee.}

Les fonctions \(f_n\) sont mesurables, puisqu'elles sont
d\'efinies sur \(\N\) muni de la tribu \(\mathcal{P}(\N)\).
Elles convergent simplement vers la fonction
\[
f(k)=e^{-\alpha k},
\]
et elles satisfont
\[
|f_n(k)|\leq g(k)
\]
pour tous \(n\geq1\) et \(k\in\N\), avec \(g\in L^1(\mu)\).

Le th\'eor\`eme de convergence domin\'ee donne donc
\[
\lim_{n\to\infty}\int_{\N}f_n\,d\mu
=
\int_{\N}
\left(\lim_{n\to\infty}f_n\right)d\mu.
\]
Par cons\'equent,
\[
\begin{aligned}
\lim_{n\to\infty}\frac{S_n}{n^{\alpha n}}
&=
\lim_{n\to\infty}
\sum_{k=0}^{\infty}f_n(k)\\
&=
\sum_{k=0}^{\infty}e^{-\alpha k}\\
&=
\frac{1}{1-e^{-\alpha}}.
\end{aligned}
\]
\end{preuvebox}

\begin{resultatbox}
\textbf{R\'esultat final.}

Pour tout r\'eel \(\alpha>0\),
\[
\boxed{
\lim_{n\to\infty}
\frac{
  1^{\alpha n}+2^{\alpha n}+\cdots+n^{\alpha n}
}{
  n^{\alpha n}
}
=
\frac{1}{1-e^{-\alpha}}
}.
\]
De mani\`ere \'equivalente,
\[
\boxed{
1^{\alpha n}+2^{\alpha n}+\cdots+n^{\alpha n}
\sim
\frac{n^{\alpha n}}{1-e^{-\alpha}}
}
\qquad
(n\longrightarrow+\infty).
\]
\end{resultatbox}

\begin{remarquebox}
\textbf{Remarque sur l'espace mesur\'e.}

Le recours \`a la mesure de comptage permet d'appliquer le
th\'eor\`eme de convergence domin\'ee sur un espace mesur\'e
ind\'ependant de \(n\). L'extension
\[
f_n(k)=0
\qquad\text{pour } k\geq n
\]
est n\'ecessaire pour que toutes les fonctions \(f_n\) soient
d\'efinies sur le m\^eme ensemble \(\N\).
\end{remarquebox}

\begin{remarquebox}
\textbf{Interpr\'etation asymptotique.}

Apr\`es normalisation par \(n^{\alpha n}\), les termes situ\'es
\`a une distance fix\'ee \(k\) de \(n\) v\'erifient
\[
\left(\frac{n-k}{n}\right)^{\alpha n}
\longrightarrow e^{-\alpha k}.
\]
La constante finale est donc la somme de toutes ces
contributions limites :
\[
1+e^{-\alpha}+e^{-2\alpha}+\cdots
=
\frac{1}{1-e^{-\alpha}}.
\]
\end{remarquebox}

\begin{remarquebox}
\textbf{Classification.}

Cet exercice appartient au th\`eme \textbf{Analyse}.
\end{remarquebox}

\end{document}